% copyright (c) 2015 Synrc Research Center

\documentclass[11pt,oneside]{article}

% Copyright (c) 2010 Synrc Research Center

\usepackage{ifthen}
\usepackage[english]{babel}
\usepackage{bussproofs}
\usepackage{tabstackengine}
\usepackage{graphicx}
\usepackage{cite}
\usepackage[usenames,dvipsnames]{color}
\usepackage[hmarginratio=2:3]{geometry}
\usepackage{inconsolata}
\usepackage{amssymb}
\usepackage{amsmath}
\usepackage[lBrack,rBrack,llparenthesis,rrparenthesis]{stmaryrd} % for comparison with 4 similar parenthesis symbols
\usepackage{accsupp}

\usepackage{mathtools}
\usepackage{hyperref}
\usepackage[utf8]{inputenc}
\usepackage[T1]{fontenc}
\usepackage{listingsutf8}
\usepackage{moreverb}
\usepackage{listings}
\usepackage[none]{hyphenat}
\usepackage{caption}
\usepackage{pgf}
\usepackage{tikz-cd}

\lstset{morekeywords={record,inductive,coinductive,extend,data}}

\hyphenation{framework nitrogen javascript facebook}

% include image for HeVeA and LaTeX

\makeatletter
\def\@seccntformat#1{\llap{\csname the#1\endcsname\quad}}
\makeatother

\newcommand{\includeimage}[2]
{\ifhevea
    {\imgsrc{#1}}
\else{
    \begin{figure}[h!]
    \centering
    \includegraphics[width=\textwidth]{#1}
    \caption{#2}
    \end{figure}}
\fi}

\lstset{
    backgroundcolor=\color{white},
    keywordstyle=\color{blue},
    basicstyle=\bf\ttfamily\footnotesize,
    columns=fixed}

\headsep = 0cm
\voffset = 0cm
\hoffset = 1cm
\topmargin = 0cm
\textwidth = 15cm
\textheight = 22cm
\footskip = 1.3cm
\parindent = 0cm

\hyphenpenalty=5000
  \tolerance=1000

\newcommand{\sign}[1]{%      
  \begin{tabular}[t]{@{}l@{}}
  \makebox[1.5in]{\dotfill}\\
  \strut#1\strut
  \end{tabular}%
}
\newcommand{\Date}{%
  \begin{tabular}[t]{@{}p{1.5in}@{}}
  \\[-1ex]
  \strut Date: \dotfill\strut
  \end{tabular}%
}

\addto\captionenglish{\renewcommand*{\bibname}{Literaturliste}}

\lstset{
 inputencoding=utf8,
extendedchars=true
}

\begin{document}

\thispagestyle{empty}
\begin{center}

\begingroup
\parbox[t][][l]{0.30\textwidth}{ \includegraphics[scale=0.3]{img/grp} }
\parbox[t][][r]{0.60\textwidth}{ \flushright \textsc{{\Large {\bf {Groupoid Infinity}}}}\\
                                             \textsc{Kostiantynivska 20/14, Kyiv, Ukraine 04071}\\  }\endgroup

\vspace{6cm}   {\Large \bf Identity Type Encoding\\}\par
\vspace{0.3cm} {\Large Paul Lyutko, Maxim Sokhatsky\par}
\vspace{6cm}   {\Large Groupoid Infinity\par}
\vspace{0.3cm} {\Large Kropyvnytsky 2016}

\end{center}

\newpage
\vspace{2cm}
\tableofcontents
\newpage

\section{Abstract}

In this article Gropoid Infinit will show how to encode classical
Identity Types used in Type Theory using EXE language. We provide brief explanation
of Identity Type properties and our motivation why use Setoid Types instead
of classical Identity Types from Type Theory.

\begin{lstlisting}[mathescape=true]
record Id (A: Type): A $\rightarrow$ A $\rightarrow$ Type :=
       (refl (a: A): Id a a)
\end{lstlisting}

\section{Identity Type}

\subsection{Typing Rules}

\begin{equation}
\tag{$Id$-formation}
\dfrac
  {a:A\ \ \ \ b:A\ \ \ \ A:Type}
  {Id(A,a,b) : Type}
\end{equation}

\begin{equation}
\tag{$Id$-intro}
\dfrac
  {a:A}
  {refl(A,a) : Id(A,a,a) }
\end{equation}

\begin{equation}
\tag{$J$-elimination}
\dfrac
  {p:Id(a,b)\ \ \ \ x,y:A\ \ \ \ u:Id(x,y) \vdash E:Type\ \ \ \ x:A \vdash d: E\ [x/y,\ refl(x)/u]}
  {J(a,b,p,(x,y,u)\ d) : E\ [a/x,\ b/y,\ p/u]}
\end{equation}

\begin{equation}
\tag{$Id$-computation}
\dfrac
  {a,x,y:A,\ \ \ \ u:Id(x,y) \vdash E:Type\ \ \ \ x:A \vdash d:E\ [x/y,\ refl(x)/u]}
  {J(a,a,refl(a),(x,y,u)\ d) = d\ [a/x] : E\ [a/y,\ refl(a)/u]}
\end{equation}

\begin{lstlisting}[mathescape=true]
record Id (A: Type): Type :=
       (Id: A $\rightarrow$ A $\rightarrow$ Type)
       (refl (a: A): Id a a)
       (Predicate := $\forall$ (x,y: A) $\rightarrow$ Id x y $\rightarrow$ Type)
       (Forall (C: Predicate) := $\forall$ (x,y: A) $\rightarrow$ $\forall$ (p: Id x y) $\rightarrow$ C x y p)
       ($\Delta$ (C: Predicate) := $\forall$ (x: A) $\rightarrow$ C x x (refl x))
       (axiom-J (C: Predicate): $\Delta$ C $\rightarrow$ Forall C)
       (computation-rule (C: Predicate) (t: $\Delta$ C):
                         $\forall$ (x: A) $\rightarrow$ (J C t x x (refl x)) ==> (t x))
\end{lstlisting}

\newpage
\subsection{Symmetry and Transitivity}

\begin{lstlisting}[mathescape=true]
record Properties (A: Type): Type :=
       (Trans (a,b,c: A) : Id a b $\rightarrow$ Id b c $\rightarrow$ Id a c :=
                           Id.axiom-J ($\lambda$ a b p1 $\rightarrow$ Id b c $\rightarrow$ Id a c) ($\lambda$ x p2 $\rightarrow$ p2) a ab)
       (Sym (a,b: A)     : Id a b $\rightarrow$ Id b a :=
                           Id.axiom-J ($\lambda$ a b p $\rightarrow$ Id b a) Id.refl a b)
\end{lstlisting}

\subsection{Substitution in Predicates}

\begin{lstlisting}[mathescape=true]
record Subst (A: Type): Type :=
       (intro (P (a: A): Type) (a1, a2: A) : Id a1 a2 $\rightarrow$ P a1 $\rightarrow$ P a2 :=
        Id.axiom-J ($\lambda$ a1 a2 p12 $\rightarrow$ P a1 $\rightarrow$ P a2) ($\lambda$ a0 p0 $\rightarrow$ p0) a1 a2)
\end{lstlisting}

\subsection{Uniqueness of Identity Proofs}

\begin{lstlisting}[mathescape=true]
record UIP (A: Type): Type :=
       (intro (A: Type) (a,b: A) (x,y: Id a b) : Id (Id A a b) x y)
\end{lstlisting}

\subsection{Axiom K}

\begin{lstlisting}[mathescape=true]
record K (A: Type): Type :=
       (Predicate$_K$ := $\forall$ (a: A) $\rightarrow$ Id a a $\rightarrow$ Type)
       (Forall$_K$ (C: Predicate$_K$) := $\forall$ (a: A) $\rightarrow$ $\forall$ (p: Id a a) $\rightarrow$ C a p)
       ($\Delta_K$ (C: Predicate$_K$) := $\forall$ (a: A) $\rightarrow$ C a (Id.refl a))
       (axiom-K (C: Predicate): $\Delta_K$ C $\rightarrow$ Forall$_K$ C)
\end{lstlisting}

\subsection{Congruence}

\begin{lstlisting}[mathescape=true]
define Respect$_1$ (A,B: Type) (C: A $\rightarrow$ Type) (D: B $\rightarrow$ Type) (R: A $\rightarrow$ B $\rightarrow$ Prop)
       (Ro: $\forall$ (x: A) (y: B) $\rightarrow$ C x $\rightarrow$ D y $\rightarrow$ Prop) :
            ($\forall$ (x: A) $\rightarrow$ C x) $\rightarrow$ ($\forall$ (x: B) $\rightarrow$ D x) $\rightarrow$ Prop :=
                       $\lambda$ (f g: Type $\rightarrow$ Type) $\rightarrow$ ($\forall$ (x y) $\rightarrow$ R x y) $\rightarrow$ Ro x y (f x) (g y)
\end{lstlisting}

\begin{lstlisting}[mathescape=true]
record Respect$_2$ (A: Type): Type :=
       (intro (A,B: Type) (f: A $\rightarrow$ B) (a,b: A) : Id A a b $\rightarrow$ Id B (f a) (f b) :=
              Id.axiom-J ($\lambda$ a b p12 $\rightarrow$ Id B (f a) (f b)) ($\lambda$ x $\rightarrow$ refl B (f x)) a b)
\end{lstlisting}

\subsection{Setoid}

\begin{lstlisting}[mathescape=true]
record Setoid: Type :=
       (Carrier: Type)
       (Equ: Carrier $\rightarrow$ Carrier $\rightarrow$ Prop)
       (Refl:  $\forall$ (e$_0$: Carrier) $\rightarrow$ Equ e$_0$ e$_0$)
       (Trans: $\forall$ (e$_1$,e$_2$,e$_3$: Carrier) $\rightarrow$ Equ e$_1$ e$_2$ $\rightarrow$ Equ e$_2$ e$_3$ $\rightarrow$ Equ e$_1$ e$_3$)
       (Sym:   $\forall$ (e$_1$ e$_2$: Carrier) $\rightarrow$ Equ e$_1$ e$_2$ $\rightarrow$ Equ e$_2$ e$_1$)
\end{lstlisting}

\subsection{Conclusion}
As you can see EXE language has enough expressive power to be used for
drawing MLTT axioms in computer science articles and papers.

\newpage
\begin{thebibliography}{9}

\bibitem{bis} E.Bishop \textit{Foundations of Constructive Analysis} 1967
\bibitem{str} T.Streicher, M.Hofmann  \textit{The groupoid interpretation of type theory} 1996
\bibitem{bar} G.Barthe, V.Capretta \textit{Setoids in type theory} 2003
\bibitem{soz} M.Sozeau, N.Tabareau \textit{Internalizing Intensional Type Theory} 2013
\bibitem{voe} V.Voevodsky \textit{Identity types in the C-systems defined by a universe category} 2015
\bibitem{con} D.Selsam and L.de Moura \textit{Congruence Closure in Intensional Type Theory} 2016
\end{thebibliography}

\end{document}
